\documentclass[12pt,a4paper]{report}
\usepackage[utf8]{inputenc}
\usepackage[brazilian]{babel}
\usepackage{geometry}
\usepackage{amsmath}
\usepackage{amssymb}
\usepackage{graphicx}
\usepackage{hyperref}
\usepackage{booktabs}
\usepackage{array}
\usepackage{float}
\usepackage{fancyhdr}\usepackage{helvet}
\usepackage{sectsty}
\usepackage{titlesec}
\usepackage{indentfirst}\usepackage{tabularx}
\usepackage{longtable}

\renewcommand{\familydefault}{\sfdefault}
\allsectionsfont{\sffamily\fontsize{12}{14}\selectfont}

\titleformat{\chapter}[hang]{\normalfont\Large\bfseries}{\thechapter.}{1em}{}
\geometry{margin=2.5cm}

\setlength{\parskip}{0.5cm}
\setlength{\parindent}{0.75cm}

\pagestyle{fancy}
\fancyhf{}
\rhead{Relatório de Reclamações - Pombal}
\lhead{\thepage}
\renewcommand{\headrulewidth}{0.4pt}

\title{\textbf{RELATÓRIO DE RECLAMAÇÕES}\\
Município de Pombal - PB\\
Questões com a Concessionária de Energia}
\author{}
\date{\today}

\begin{document}

\maketitle

\tableofcontents
\newpage

\chapter{Objetivo}

O presente relatório tem por finalidade apresentar um panorama atualizado das reclamações administrativas em tramitação relacionadas às unidades consumidoras vinculadas ao município de \textbf{Pombal-PB}, indicando o assunto tratado, o objeto da demanda, a unidade consumidora envolvida e o atual status de análise. 

As reclamações encontram-se em acompanhamento perante a Ouvidoria da ANEEL e a Ouvidoria da concessionária, tratando principalmente de descumprimento de cobranças indevidas e pedidos de devolução complementar. 

\vspace{3cm}

\chapter{Resumo Geral das Reclamações}

Atualmente, foram identificadas 22 reclamações, a maioria concluída e algumas ainda em análise, como mostra a tabela \ref{tabela_resumo}.

\setlength{\extrarowheight}{0.5cm}
\begin{longtable}{|>{\centering\arraybackslash}m{8cm}|>{\centering\arraybackslash}m{8cm}|}
\caption{Resumo de Reclamações por Situação} \label{tabela_resumo} \\
\hline
\textbf{Situação} & \textbf{Quantidade} \\
\hline
\endfirsthead
\hline
\textbf{Situação} & \textbf{Quantidade} \\
\hline
\endhead
\hline
\endfoot
Concluída & 15 \\
\hline
Em análise na ouvidoria da ANEEL & 1 \\
\hline
Pendente contestação & 6 \\
\hline
\textbf{Total} & \textbf{22} \\
\hline
\end{longtable}

A maior parte das demandas ainda em tramitação estão com pendencia de contestação na ouvidoria da concessionária. 

\vspace{3cm}

\chapter{Principais Teses em Discussão}

A tese de maior recorrência é a de cobrança indevida, principallemnte em função de  enquadramento AES, presente em 12 reclamações,e Enquadrammento IP, presente em 6 reclamações.  

Também há volume relevante de reclamações de outros tipos de cobrança indevida, como parcela DIC/2, valor incompatível com QIP e perda nos reatores.

\setlength{\extrarowheight}{0.5cm}
\begin{longtable}{|>{\centering\arraybackslash}m{8cm}|>{\centering\arraybackslash}m{8cm}|}
\caption{Resumo de teses em discussão} \label{tabela_teses} \\
\hline
\textbf{Tese/assunto} & \textbf{Quantidade} \\
\hline
\endfirsthead
\hline
\textbf{Tese/assunto} & \textbf{Quantidade} \\
\hline
\endhead
\hline
\endfoot
Enquadramento IP & 6 \\
\hline
Enquadramento AES & 12 \\
\hline
Não isenção de imposto de renda & 1 \\
\hline
Parcela DIC/2 & 1 \\
\hline
Valor incompativel com QIP & 1 \\
\hline
Perda nos reatores & 1 \\
\hline
\textbf{Total} & \textbf{22} \\
\hline
\end{longtable}

\vspace{3cm}

\chapter{Planilhas Detalhadas de Reclamações}

\setlength{\extrarowheight}{0.5cm}
\begin{longtable}{|>{\centering\arraybackslash}m{3cm}|>{\centering\arraybackslash}m{3cm}|>{\centering\arraybackslash}m{5cm}|>{\centering\arraybackslash}m{4cm}|}
\caption{Detalhamento das Reclamações} \\
\hline
\textbf{Reclamação} & \textbf{Unidade Consumidora} & \textbf{Objeto da demanda} & \textbf{Status} \\
\hline
\endfirsthead
\hline
\textbf{Reclamação} & \textbf{Unidade Consumidora} & \textbf{Objeto da demanda} & \textbf{Status} \\
\hline
\endhead
\hline
\endfoot
002/2025 & 385967-5 & Enquadramento AES & Pendente de contestação \\
\hline
003/2025 & 1389500-8 & Enquadramento AES & Concluído \\
\hline
004/2025 & 1389600-6 & Enquadramento AES & Concluído \\
\hline
005/2025 & 1389624-6 & Enquadramento AES & Concluído \\
\hline
006/2025 & 1389638-6 & Enquadramento AES & Pendente de contestação \\
\hline
007/2025 & 1389646-9 & Enquadramento AES & Concluído \\
\hline
008/2025 & 1389869-7 & Enquadramento AES & Concluído \\
\hline
009/2025 & 1390173-1 & Enquadramento AES & Concluído \\
\hline
010/2025 & 1390226-7 & Enquadramento AES & Pendente de contestação \\
\hline
011/2025 & 1571531-1 & Enquadramento AES & Concluído \\
\hline
012/2025 & 1599798-4 & Enquadramento AES & Pendente de contestação \\
\hline
013/2025 & 1882785-7 & Enquadramento AES & Concluído \\
\hline
014/2025 & 12630-0 & Enquadramento IP & Concluído \\
\hline
015/2025 & 536551-5 & Enquadramento IP & Concluído \\
\hline
016/2025 & 1066648-5 & Enquadramento IP & Concluído \\
\hline
017/2025 & 1271871-4 & Enquadramento IP & Concluído \\
\hline
018/2025 & 1452856-6 & Enquadramento IP & Concluído \\
\hline
019/2025 & 1453356-6 & Enquadramento IP & Concluído \\
\hline
020/2025 & 207926 & Perdas nos reatores & Em análise na ouvidoria da ANEEL \\
\hline
021/2025 & S/I & Não isenção do Imposto de Renda & Concluído \\
\hline
022/2025 & 207926 & Ausência parcela DIC/2 & Concluído \\
\hline
001/2026 & 207926 & IP Estimada & Pendente de contestação \\
\hline
\end{longtable}

\setlength{\extrarowheight}{0.5cm}
\begin{longtable}{|>{\centering\arraybackslash}m{3cm}|>{\centering\arraybackslash}m{7cm}|>{\centering\arraybackslash}m{5cm}|}
\caption{Resumo Financeiro das Teses Deferidas} \label{tabela_resumo_financeiro} \\
\hline
\textbf{Reclamação} & \textbf{Tese Deferida} & \textbf{Valor Deferido (R\$)} \\
\hline
\endfirsthead
\hline
\textbf{Reclamação} & \textbf{Tese Deferida} & \textbf{Valor Deferido (R\$)} \\
\hline
\endhead
\hline
\endfoot
005/2025 & Enquadramento AES & 686,56 \\
\hline
008/2025 & Enquadramento AES & 3.561,71 \\
\hline
011/2025 & Enquadramento AES & 6.811,96 \\
\hline
020/2025 & Perdas nos reatores & 213.668,79 \\
\hline
022/2025 & Ausencia de parcela DIC/2 & 7.797,86 \\
\hline
\textbf{Total} &  & \textbf{232.526,88} \\
\hline
\end{longtable}



\vspace{3cm}

\chapter{Análise Gerencial}

As demandas de Pombal estão concentradas em temas de impacto financeiro. 

O primeiro ponto de atenção refere-se aos enquadramentos incorretos, que representam a maior parte das reclamações em andamento. 

O segundo ponto relevante envolve as reclamações que ainda serão contestadas, essas poderão gerar devoluções significativas, principalmente as relacionadas ao enquadramento AES, que tem potencial de gerar devolução de aproximadamente 10 mil.

Há ainda reclamação específica sobre solicitação de devolução complementar referente ao período de 10 anos, tema que deve gerar retorno financeiro de aproximadamente 60 mil. 
\section*{}

\vspace{3cm}

\chapter{Conclusão}


\end{document}
